\renewcommand\epsilon{\varepsilon}
\renewcommand\hat{\widehat}
\newcommand\proba{\frac{1}{\size \F}\left(p_i + \frac{4}{\epsilon} - 1\right)\left(\frac{4}{\epsilon}\right)^{p_i}}
\newcommand\F{\mathbb{F}}
\newcommand\Fq{\mathbb{F}_q}
\newcommand\PP{\mathbb{P}}
\newcommand\Z{\mathbb{Z}}
\newcommand\N{\mathbb{N}}

\newcommand\floor[1]{\left\lfloor #1 \right\rfloor}
\newcommand\ceil[1]{\left\lceil #1 \right\rceil}
\newcommand\gen[1]{\left\langle #1 \right\rangle}
\newcommand\fold[1]{\textsf{\textbf{Fold}}\left[#1\right]}
\newcommand\poles[1]{\div_\infty\left( #1 \right)}
\newcommand\zeroes[1]{\div_0\left( #1 \right)}

\renewcommand{\hat}{\widehat}
\renewcommand{\bar}{\overline}

\DeclareMathOperator{\polylog}{polylog}
\DeclareMathOperator{\Supp}{Supp}
\DeclareMathOperator{\Aut}{Aut}
\DeclareMathOperator{\Div}{Div}
\let\div\relax\DeclareMathOperator{\div}{div}

\newcommand\Inv[2]{\operatorname{Inv}_{#1}(#2)}

\def\rest#1#2{{{#1}_{|#2}}}

\newcommand\RS[1]{\mathsf{RS}\left[\F,#1\right]}

\newcommand\prover{\mathsf{P}}
\newcommand\verifier{\mathsf{V}}

\renewcommand{\vec}[1]{\boldsymbol{#1}}
\newcommand{\set}[1]{\left\{#1\right\}}
\newcommand{\size}[1]{\left|#1\right|}
\newcommand{\range}[1]{\set{0,\dots,#1-1}}
\newcommand{\Range}[2]{\set{#1,\dots,#2}}

\newcommand\bfz{{\vec{z}}}
\newcommand\bfzero{{\vec{0}}}

\newcommand\calY{\mathcal{Y}}
\newcommand\calX{\mathcal{X}}
\newcommand\calL{\mathcal{L}}
\newcommand\calH{\mathcal{H}}
\newcommand\calF{\mathcal{F}}
\newcommand\calP{\mathcal{P}}
\newcommand\calI{\mathcal{I}}
\newcommand\calT{\mathcal{T}}
\newcommand\calC{\mathcal{C}}
\newcommand\calG{\mathcal{G}}

\newcommand\seqG{\vec{\calG}}
\newcommand\XGseq{(\calX, \vec{\calG})}

\newcommand\xpi[1]{#1^{(i)}}
\newcommand\xpip[1]{#1^{(i+1)}}

\newcommand{\err}{\textsf{err}}
\newcommand{\la}{\langle}
\newcommand{\ra}{\rangle}

\newcommand\mydef{\coloneqq}

\newcommand{\cd}{\cdot}
\newcommand{\ii}{i_{\rm{max}}}

\providecommand{\proofsubparagraph}{\paragraph}
\makeatletter
\@ifundefined{notation}{\newtheorem{notation}{Notation}}{}
\makeatother

% Restore the LIPIcs autoref alias counters and printed type names.
\makeatletter
\theoremstyle{plain}
\@ifundefined{theorem}{\newtheorem{theorem}{Theorem}}{}
\theoremstyle{plain}
\@ifundefined{lemma}{\newaliascnt{lemma}{theorem}\newtheorem{lemma}[lemma]{Lemma}\aliascntresetthe{lemma}}{}
\providecommand{\lemmaautorefname}{Lemma}
\@ifundefined{corollary}{\newaliascnt{corollary}{theorem}\newtheorem{corollary}[corollary]{Corollary}\aliascntresetthe{corollary}}{}
\providecommand{\corollaryautorefname}{Corollary}
\@ifundefined{proposition}{\newaliascnt{proposition}{theorem}\newtheorem{proposition}[proposition]{Proposition}\aliascntresetthe{proposition}}{}
\providecommand{\propositionautorefname}{Proposition}
\@ifundefined{definition}{\newaliascnt{definition}{theorem}\newtheorem{definition}[definition]{Definition}\aliascntresetthe{definition}}{}
\providecommand{\definitionautorefname}{Definition}
\theoremstyle{definition}
\@ifundefined{example}{\newaliascnt{example}{theorem}\newtheorem{example}[example]{Example}\aliascntresetthe{example}}{}
\providecommand{\exampleautorefname}{Example}
\theoremstyle{remark}
\@ifundefined{remark}{\newaliascnt{remark}{theorem}\newtheorem{remark}[remark]{Remark}\aliascntresetthe{remark}}{}
\providecommand{\remarkautorefname}{Remark}
\providecommand{\theoremautorefname}{Theorem}
\providecommand{\notationautorefname}{Notation}
\renewcommand{\sectionautorefname}{Section}
\renewcommand{\subsectionautorefname}{Subsection}
\renewcommand{\subsubsectionautorefname}{Subsubsection}
\theoremstyle{plain}
\makeatother
